\section{Limitations}
\label{sec:limitations}

\begin{enumerate}
  \item \textbf{Scope of the identity.} Lemma~\ref{lem:exact-stage-identity} is stated for the open
  chain, where the joint coordinates form a global chart. For closed loops the same rows are solved,
  but the identity holds only on local charts of independent coordinates, and the theorem does not
  by itself turn the roundoff-level constraint residuals observed on the four-link and slider-crank
  mechanisms into a trajectory-error bound; that needs a stability constant for the mechanism at
  hand.
  \item \textbf{Regular branch.} The frozen sliding direction of the chain's second pair is a
  modelling choice that is regular only while $n_1\cdot a_1(q)\neq0$; on the benchmark initial state
  this holds up to $t\approx0.6$, and all chain experiments stop at $T=0.5$. A pair whose sliding
  direction follows the axis would remove the limit at the cost of a different (still smooth) stage
  system; the analysis applies unchanged.
  \item \textbf{Sharp friction.} The theorem's constants and its threshold $h_0$ are governed by the
  curvature of the regularized friction law. For Stribeck velocities below about $0.1$ the observed
  order is reduced on coarse grids and recovers only for $h\lesssim3\times10^{-3}$
  (Section~\ref{sec:numerical}). Set-valued Coulomb friction is outside the smoothness hypothesis.
  \item \textbf{Solver hypothesis.} Lemma~\ref{lem:newton-envelope} covers predictors within $O(h)$
  of the Gauss stage. The implemented predictor is not of that kind, and the experiments use a fixed
  Newton tolerance rather than the $h^7$-scaled rule, so on the finest grids they satisfy (H3) only
  with a large constant $c_\eta$; that the accepted iterates are nevertheless the Gauss roots to
  within $10^{-11}$ is checked a posteriori, not guaranteed a priori.
  \item \textbf{Cost.} The reference implementation forms and factorizes a dense
  $132\times132$ Jacobian by automatic differentiation in Python; wall-clock times are dominated by
  interpreter overhead and are reported only as indications. A production implementation would
  exploit the stage-block structure and the sparsity of the constraint rows; no scalability study
  in the number of bodies is reported.
  \item \textbf{Comparisons.} The time finite-element method of~\cite{chaturvedi2026tfe} has not
  been implemented; the comparison with it is formal (order six against order five). The public
  absolute-coordinate codes are run as published, with their own tolerances, and are compared on the
  double pendulum only, the one benchmark mechanism whose motion is not prescribed.
\end{enumerate}

\section{Conclusions}
\label{sec:conclusions}

The \method{} step solves a single square stage system that couples three-stage Gauss collocation
of the joint coordinates, the lower-pair constraints at position, velocity and acceleration level,
and the Newton--Euler balance, and reconstructs the rotational endpoint through the exponential
map. The stage system is exactly reduced Gauss collocation lifted to absolute coordinates: the
lifted reduced Gauss stage satisfies all $132$ rows, and every root of the non-dynamic rows on the
regular branch is such a lift. The sixth order of the method is therefore inherited from Gauss
collocation, and the implementation-side perturbations, endpoint closure and inexact Newton, are
$O(h^7)$ under a smoothness hypothesis alone; the interfaces through which the implementation enters
are derived for small steps, with a threshold that is set by the friction law. The algebra and the
perturbation chain are machine-checked.

The experiments confirm the picture at every point where it can be tested: sixth order in position,
orientation, velocity and angular velocity over the regular branch of a frictional chain, with
constraint norms at the $10^{-13}$ level and energy conservation of the frictionless variant to
$10^{-12}$; orders two, four and six for the one-, two- and three-stage members with the same
stage system at the same Newton work; sixth order on the ASME double pendulum and roundoff-level
reproduction of the three driven mechanisms; and a controlled loss of order when the friction law is
sharpened, in the direction the analysis predicts. Against the public absolute-coordinate codes on
the double pendulum over the public horizon, the method reaches errors below $10^{-10}$ at step
sizes where those codes still have errors of order one.

Three directions follow. The identity suggests that other stage-level constraint hierarchies
(Lobatto, Radau) inherit their reduced-problem order in the same way, which would give
$L$-stable members for stiff joints. The $O(h)$ predictor of Lemma~\ref{lem:newton-envelope}
would remove the last a-posteriori element of the analysis at no cost. And a sparse, block-structured
implementation would turn the per-step cost into one comparable to that of the low-order
absolute-coordinate schemes, at which point the work/precision advantage measured here in Newton
iterations would appear in wall-clock time as well.
